% 337_Zeit_Emergenz_FFGFT_GALG_De_ch.tex
% Johann Pascher, 28. Aug. 2026

\providecommand{\BM}{\textbf{[B]}}
\providecommand{\KM}{\textbf{[K]}}
\providecommand{\SM}{\textbf{[S]}}
\providecommand{\SetzM}{\textbf{[SETZUNG]}}

\chapter*{Zeit-Emergenz und Gravitationskonstante: Korrekturvorschläge zu Matzkes GALG-Rahmen}
\addcontentsline{toc}{chapter}{Korrekturvorschläge zu Matzkes GALG-Rahmen}

\begin{abstract}
Dieses Dokument formuliert zwei Korrekturvorschläge zum GALG-Framework
von Doug Matzke~\cite{Matzke2002,Matzke2026IPI}, gestützt auf direkte
Zitate und algebraische Nachrechnung.

\textbf{Vorschlag~1 (Zeit):}
Matzkes Co-exclusion-Primitiv erzeugt proto-Zeit durch
Operator-Wechsel~\cite{Matzke2002}~(S.~37), analog zur FFGFT-Dekompaktifizierung
(Ausrollen, Dok.~333). Beide Frameworks haben dieselbe drei-schichtige
Emergenz-Struktur. Das ist keine Schwäche von GALG, sondern eine bisher
nicht explizit gemachte Parallele, die Dougs Rahmen stärkt.

\textbf{Vorschlag~2 (G):}
Die in~\cite{Matzke2026IPI} behauptete ``Ableitung'' der Schwarzschild-Geometrie
``ohne freie Parameter'' verwendet $M_\text{bit} = m_P\cdot\ln2/(2\pi)$
als Startpunkt. Da $m_P = \sqrt{\hbar c/G}$ per Definition, steckt~$G$
bereits im Eingang. Das implizite~$G$ stimmt numerisch exakt mit~$G_\text{SI}$
überein (0{,}00\,\%) --- aber als Tautologie, nicht als Ableitung.
Der Vorschlag: diese Einschränkung explizit benennen,
und den Rahmen für eine echte $G$-Ableitung offen halten.
In FFGFT ist $G = \xi^2\hbar c/(4m_e^2)\cdot K_\text{frak}$ ohne~$m_P$
als Eingang abgeleitet (Dok.~012).

Beide Vorschläge werden durch direkte Zitate aus Matzkes Veröffentlichungen
und durch Prüfskripte belegt.
\end{abstract}

% ============================================================
\section*{Leitfaden}
\addcontentsline{toc}{section}{Leitfaden}
% ============================================================

\textbf{Das Grundproblem.}
Beobachtbare Zeit ist ein makroskopisches Phänomen. Quantenmechanik und
algebraische Feldtheorien arbeiten auf einer Ebene, wo Zeit entweder
fehlt oder nur als diskretes Primitiv erscheint.
Wie entsteht die kontinuierliche, erlebbare Zeit aus einer zeitlosen
algebraischen Grundstruktur?

\textbf{Was dieses Dokument zeigt.}
Zwei präzise Korrekturvorschläge an Doug Matzke, gestützt auf
direkte Zitate aus seinen eigenen Publikationen:
(1)~Die Zeit-Emergenz in GALG hat dieselbe Struktur wie in FFGFT ---
das stärkt Dougs Rahmen, ist aber noch nicht explizit gemacht.
(2)~Die behauptete G-Ableitung in~\cite{Matzke2026IPI} ist eine Tautologie
--- das begrenzt den Anspruch, schwächt aber nicht die Algebra.

% ============================================================
\section*{Zeichenerklärung}
\addcontentsline{toc}{section}{Zeichenerklärung}
% ============================================================

\begin{center}
\begin{tabular}{lll}
\toprule
\textbf{Symbol} & \textbf{Bedeutung} & \textbf{Quelle} \\
\midrule
\multicolumn{3}{l}{\textit{Statusmarker}} \\
\BM  & algebraisch bewiesen, Prüfskript vorhanden & Dok.~190 \\
\KM  & aus $\xi$ ableitbar & Dok.~190 \\
\SM  & offene Brücke, Skizze & Dok.~190 \\
\midrule
\multicolumn{3}{l}{\textit{FFGFT}} \\
$\tilde{T} = 1/m$ & intrinsische Zeitskala einer Mode & Dok.~321~\cite{Pascher321} \\
$T^4/\mathbb{Z}_3$ & kompakte 4-Torus-Geometrie & Dok.~257~\cite{Pascher257} \\
Dekompaktifizierung & Ausrollen der kompakten Dimension & Dok.~333~\cite{Pascher333} \\
\midrule
\multicolumn{3}{l}{\textit{Matzke GALG}} \\
Co-occurrence & Addition $a+b$: räumliches Primitiv & \cite{Matzke2002} S.~37 \\
Co-exclusion & Multiplikation $ab$: zeitliches Primitiv & \cite{Matzke2002} S.~37 \\
Proto-Zeit & emergente Zeitordnung aus Co-exclusion & \cite{Matzke2002} S.~14 \\
\bottomrule
\end{tabular}
\end{center}

% ============================================================
\section{Zeit in FFGFT: kompaktifiziert und emergent}
% ============================================================

\subsection{Die kompaktifizierte Zeitskala $\tilde{T} = 1/m$}

In FFGFT ist jedem Teilchen der Masse $m$ eine intrinsische
Zeitskala zugeordnet (Dok.~321~\cite{Pascher321}):
\begin{equation}
  \tilde{T} = \frac{1}{m}
  \qquad \text{(natürliche Einheiten, $\hbar = c = 1$)}
\end{equation}
Diese Zeitskala lebt auf der kompakten Geometrie T$^4/\mathbb{Z}_3$.
Sie ist modengebunden --- jede Mode hat ihre eigene Zeitskala,
und modenübergreifende Superpositionen erzeugen inkompatible
Zeitskalen (Dok.~334~\cite{Pascher334}) \BM.

\subsection{Dekompaktifizierung als Emergenz-Mechanismus}

Die kompakte Dimension ist nicht direkt beobachtbar.
Beobachtbare Zeit entsteht erst beim Ausrollen:
dem Übergang von der kompakten Geometrie in die makroskopische
Raumzeit (Dok.~333~\cite{Pascher333}).

Dok.~333 unterscheidet zwei Ausroll-Typen:
\begin{itemize}
  \item \textbf{Typ~I} (Messakt): verlustbehaftet, irreversibel ---
    der Kollaps der Wellenfunktion als Dekompaktifizierung.
  \item \textbf{Typ~III} (Pullback): verlustfrei, reversibel ---
    geometrischer Rückzug auf die kompakte Mannigfaltigkeit.
\end{itemize}
Beobachtbare Zeit ist ein Typ-I-Phänomen: sie entsteht durch
den irreversiblen Ausroll-Akt, nicht durch die kompakte Geometrie selbst.

\textbf{Kernaussage FFGFT} \BM:
Vor der Dekompaktifizierung existiert nur $\tilde{T} = 1/m$
als modengebundene geometrische Größe.
Klassische, messbare Zeit $t$ ist das Ergebnis des Ausrollens.

% ============================================================
\section{Zeit in GALG: proto-temporal und emergent}
% ============================================================

\subsection{Co-occurrence und Co-exclusion}

Matzke~\cite{Matzke2002} definiert zwei fundamentale Primitive
für die geometrische Algebra (Kap.~4.2, S.~36--39):

\textbf{Co-occurrence} (räumliches Primitiv): Zwei Zustände
existieren gleichzeitig, ausgedrückt als GA-Addition $a+b$.
Wörtlich~\cite{Matzke2002}, S.~37:
\begin{quote}
``Co-occurrence is defined as two or more states `occurring'
simultaneously [...] commutative properties of addition reflects
the lack of temporal ordering characteristic of simultaneity.''
\end{quote}
Co-occurrence ist statisch --- keine Zeitordnung.

\textbf{Co-exclusion} (zeitliches Primitiv): Ein Operator
wechselt den Zustand durch GA-Multiplikation $ab$.
Wörtlich~\cite{Matzke2002}, S.~37:
\begin{quote}
``Since co-exclusion infers an operator for change and implies
temporal sequence, it is considered to be the temporal primitive.
Conversely, co-occurrence is labeled the spatial primitive
since it is static.''
\end{quote}

\subsection{Proto-Zeit emergiert aus Operatoren}

Klassische Zeit ist nicht im GALG-Framework kodiert,
sondern entsteht aus dem Zusammenspiel disjunkter Operatoren.
Wörtlich~\cite{Matzke2002}, S.~14:
\begin{quote}
``Interaction of disjoint finite sets of dimensions
(operators and resulting co-exclusions) creates a
synchronization-based proto-time from which classical time
and energy metrics ultimately emerge.''
\end{quote}

Und zur Notwendigkeit eines anderen Zeitbegriffs im Quantenbereich
(S.~12):
\begin{quote}
``The quantum domain must have a different kind of primitive
temporal framework, if for no other reason than it exists in
a completely different spatial environment. Conceptually, this
proto-temporal framework is identical to the synchronization
tokens and mechanisms used by asynchronous logic.''
\end{quote}

\textbf{Kernaussage GALG}:
Vor der Emergenz existiert nur das algebraische Primitiv
der Co-exclusion (Multiplikation).
Klassische Zeit und Energie entstehen aus der Interaktion
disjunkter Operator-Mengen.

% ============================================================
\section{Die strukturelle Parallele}
% ============================================================

\begin{center}
\begin{tabular}{lll}
\toprule
\textbf{Schicht} & \textbf{FFGFT} & \textbf{GALG} \\
\midrule
Zeitlose Grundschicht
  & $\tilde{T}=1/m$ kompaktifiziert
  & Co-occurrence (Addition) \\
Zeitliches Primitiv
  & Torsion / Wicklung auf T$^4$
  & Co-exclusion (Multiplikation $ab$) \\
Emergenz-Mechanismus
  & Dekompaktifizierung (Ausrollen)
  & Interaktion disjunkter Operatoren \\
Ergebnis
  & beobachtbare Zeit $t$, modengebunden
  & klassische Zeit + Energie \\
\bottomrule
\end{tabular}
\end{center}

Beide Frameworks teilen dieselbe dreistufige Struktur:
(1)~zeitlose algebraische Grundschicht,
(2)~ein zeitliches Primitiv das Sequenz impliziert ohne selbst
kontinuierlich zu sein,
(3)~ein Emergenz-Mechanismus der klassische, messbare Zeit erzeugt.

% ============================================================
\section{Offene Frage}
% ============================================================

Ist Dougs Co-exclusion der algebraische Ausdruck
derselben Struktur die FFGFT als Dekompaktifizierung beschreibt?

Konkret: In FFGFT ist der Ausroll-Typ~I (Messakt, Kollaps) ein
irreversibler Übergang --- genau wie Co-exclusion einen
Zustandswechsel durch einen Operator beschreibt, ohne Möglichkeit
der Umkehr (Matzke~\cite{Matzke2002}, S.~145:
\textit{``The process is irreversible because a multiplicative
inverse does not exist for $(1\pm a_0)$ and $(1\pm a_1)$''}).

Diese Parallele --- Irreversibilität als gemeinsames Merkmal ---
könnte der algebraische Kern der Verbindung sein~\SM.

Eine Prüfung würde erfordern:
\begin{enumerate}
  \item Explizite Konstruktion des Ausroll-Operators in GALG-Notation.
  \item Nachweis dass Co-exclusion + Irreduzibilität
    die FFGFT-Massenformel $\tilde{T}=1/m$ reproduziert.
\end{enumerate}

% ============================================================
\section{Masse und Gravitationskonstante: abgeleitet oder implizit?}
% ============================================================

\subsection{Dougs Schwarzschild-Herleitung und das implizite G}

In seinem IPI-Letters-Paper 2026~\cite{Matzke2026IPI} behauptet Doug:
\begin{quote}
``This paper derives the complete black hole picture ---
Schwarzschild radius, Bekenstein entropy, Hawking radiation,
collapse threshold, and information conservation --- from a single
algebraic axiom with no general relativity and no free parameters.''
\end{quote}

Die Herleitung verwendet als Startpunkt
$M_\text{bit} = m_P \cdot \ln 2 / (2\pi)$
(Landauer-Prinzip bei $T_\text{Planck}$, dann $E = mc^2$).
Aus der Kollaps-Bedingung Compton = Schwarzschild folgt:
\begin{equation}
  \frac{\hbar}{M_\text{coll}\,c} = \frac{2G\,M_\text{coll}}{c^2}
  \quad\Rightarrow\quad
  G = \frac{\hbar c}{2\,M_\text{coll}^2}
\end{equation}
mit $M_\text{coll} = n_\text{thresh}\cdot M_\text{bit} = m_P/\sqrt{2}$
(numerisch verifiziert: Übereinstimmung 0{,}00\,\%).

Das implizite G aus dieser Konstruktion:
\begin{equation}
  G_\text{doug} = \frac{\hbar c}{2\,M_\text{coll}^2}
  = \frac{\hbar c}{m_P^2}
  = 6{,}6743\times10^{-11}\;\text{m}^3/(\text{kg}\cdot\text{s}^2)
\end{equation}
Abweichung vom experimentellen Wert: $0{,}00\,\%$.

\subsection{Warum das keine echte G-Ableitung ist}

Der Startpunkt $M_\text{bit} = m_P \cdot \ln2/(2\pi)$ enthält $m_P$,
und die Planck-Masse ist definiert als $m_P = \sqrt{\hbar c / G}$.
Das bedeutet:
\begin{equation}
  G = \frac{\hbar c}{m_P^2}
\end{equation}
ist eine \textbf{Tautologie}: $G$ steckt in $m_P$, geht über $M_\text{bit}$
in $M_\text{coll}$ ein, und wird dann wieder herausgelöst.
Das Ergebnis ist algebraisch konsistent~\SM, aber keine unabhängige
Ableitung von~$G$.

Doug selbst bestätigt dies implizit in seiner IPI-Mail vom 27.~Aug.~2026:
\begin{quote}
``the fact that $M_\text{bit}$ is directly related to $M_\text{Planck}$
was a happy coincident based on Planck units consolidation only.''
\end{quote}

\subsection{FFGFT: G als echte Ableitung}

In FFGFT wird $G$ aus $\xi$ und $m_e$ abgeleitet (Dok.~012~\cite{Pascher012}),
\emph{ohne} $m_P$ als Eingangsparameter:
\begin{equation}
  G = \frac{\xi^2\,\hbar c}{4\,m_e^2}\cdot K_\text{frak}
  \qquad \text{mit } \xi = \frac{4}{30\,000},\; K_\text{frak} = \frac{74}{75}
\end{equation}
Hier ist $\xi$ aus der T$^4/\mathbb{Z}_3$-Geometrie und $m_e$ aus der
Massenformel unabhängig von $G$ bestimmt.
Das ist eine echte Ableitung: $G$ ist Ergebnis, nicht Voraussetzung.

\subsection{Vergleich}

\begin{center}
\begin{tabular}{lll}
\toprule
\textbf{Größe} & \textbf{GALG (Matzke)} & \textbf{FFGFT} \\
\midrule
Eingang      & $m_P$ (enthält $G$!) & $\xi$, $m_e$ \\
Masse        & $M_\text{bit} = m_P\cdot\ln2/(2\pi)$ & $\tilde{T}=1/m$, dann $\xi$ \\
Gravitation  & $G = \hbar c/m_P^2$ (Tautologie) & $G = \xi^2\hbar c/(4m_e^2)\cdot K$ \\
Status       & algebraisch konsistent \SM & echte Ableitung \KM \\
\bottomrule
\end{tabular}
\end{center}

Beide Frameworks erhalten denselben numerischen Wert für $G$,
aber auf fundamental verschiedenen Wegen:
GALG durch Konsistenz mit der Planck-Definition,
FFGFT durch unabhängige Ableitung aus geometrischen Parametern.

% ============================================================
\section*{Vermerk zur Verwendung von KI-Assistenz}
\addcontentsline{toc}{section}{KI-Vermerk}
% ============================================================

\noindent
Dieses Dokument wurde unter Verwendung von KI-gestützter Assistenz
(Claude, Anthropic) erstellt.
Alle direkten Zitate wurden aus Matzkes PhD~\cite{Matzke2002}
entnommen und durch Seitenangaben belegt.
Die inhaltliche Verantwortung liegt beim Autor.
Methodische Grundsätze: A-Serie, Dok.~190.

% ============================================================
\section*{Literatur}
\addcontentsline{toc}{section}{Literatur}
% ============================================================

\begin{thebibliography}{99}

\bibitem{Matzke2002}
D.~J. Matzke,
\textit{Quantum Computation Using Geometric Algebra},
Ph.D.-Dissertation, University of Texas at Dallas, 2002.
\url{https://www.matzkefamily.net/doug/papers/PHD/phd_matzke_final.pdf}\\
\textit{Direkte Zitate in diesem Dok.: S.~12 (proto-temporales Framework),
S.~14 (proto-Zeit aus Co-exclusion), S.~37 (zeitliches Primitiv),
S.~145 (Irreversibilität).}

\bibitem{Matzke2025IPI}
D.~J. Matzke,
\textit{GALG: Geometric Algebra over $\mathbb{Z}_3\mathbb{C}$},
IPI-Vortrag, 26.~Juli 2025.
\url{https://informationphysicsinstitute.blogspot.com/2025/07/ipi-talk-doug-matzke-wwwquantumdougcom.html}

\bibitem{Pascher257}
J.~Pascher,
\textit{Dok.~257: T$^4$/$\mathbb{Z}_3$-Geometrie und FFGFT},
FFGFT-Korpus, ORCID 0009-0000-6518-4064.
\url{https://github.com/jpascher/T0-Time-Mass-Duality}

\bibitem{Pascher321}
J.~Pascher,
\textit{Dok.~321: SU(3) und $\mathbb{Z}_3$-Emergenz;
intrinsische Zeitskala $\tilde{T}=1/m$},
FFGFT-Korpus, ORCID 0009-0000-6518-4064.

\bibitem{Pascher333}
J.~Pascher,
\textit{Dok.~333: $K_\text{frak}$, Ausrollen und die
$\tilde{T}\cdot m$-Dualität},
FFGFT-Korpus, 27.~August 2026, ORCID 0009-0000-6518-4064.

\bibitem{Pascher334}
J.~Pascher,
\textit{Dok.~334: Superposition ohne Zeit},
FFGFT-Korpus, 27.~August 2026, ORCID 0009-0000-6518-4064.


\bibitem{Matzke2026IPI}
D.~J. Matzke,
\textit{Black Holes from Hyperbits: The Event Horizon as Nilpotent Algebra,
Planck Voxels, and High-Dimensional Corner Geometry},
IPI Letters 2026, Vol.~4(3): N1--N22.
\url{https://ipipublishing.org/index.php/ipil/article/download/378/201/1385}

\bibitem{Pascher012}
J.~Pascher,
\textit{Dok.~012: Herleitung der Gravitationskonstante $G$ aus $\xi$},
FFGFT-Korpus, ORCID 0009-0000-6518-4064.
\url{https://github.com/jpascher/T0-Time-Mass-Duality}

\end{thebibliography}
